\documentclass[12pt]{article}
\usepackage[utf8]{inputenc}
\usepackage[T1]{fontenc}
\usepackage[margin=1in]{geometry}
\usepackage{setspace}
\usepackage{xcolor}
\usepackage[colorlinks=true, urlcolor=blue!60!black, linkcolor=blue!60!black]{hyperref}
\setlength{\parindent}{2em}
\setlength{\parskip}{1.25em}
\renewcommand{\baselinestretch}{1.0}
\begin{document}
\begin{center}
{\LARGE \textit{Tinker v. Des Moines Independent Community School District}, 393 U.S. 503 (1969)\par}
\vspace{0.8em}
{\large Case Brief --- Four-Page Edition\par}
\end{center}
\vspace{1em}

\section{Memory Jogger}

Holding 7--2 that students do not ``shed their constitutional rights to freedom of speech or expression at the schoolhouse gate,'' the Court ruled that a school may not punish passive political expression, such as black armbands protesting the Vietnam War, absent actual or reasonably forecast material and substantial disruption of schoolwork or discipline, or an invasion of the rights of others (\href{https://supreme.justia.com/cases/federal/us/393/503/}{393 U.S. 503}).

\section{Detailed Case Facts}

In December 1965, a group of adults and students in Des Moines, Iowa, meeting at the Eckhardt home, decided to protest the Vietnam War and support a Christmas truce by wearing black armbands during the holidays and fasting on December 16 and New Year's Eve. Among them were Mary Beth Tinker (13), her brother John Tinker (15), and Christopher Eckhardt (16). Learning of the plan, the city's principals met on December 14 and adopted a policy that any student wearing an armband would be asked to remove it and suspended on refusal. The district did not ban other political symbols: students wore campaign buttons and even the Iron Cross without penalty.

Mary Beth and Christopher wore armbands on December 16, and John the next day. All three refused to remove them and were suspended until after New Year's Day. Of 18,000 students, only a handful wore armbands. A few students made hostile remarks outside class, but there were no threats, violence, or disruption of schoolwork.

\section{Procedural History}

Through their fathers, the students sued under \S~1983 for an injunction and nominal damages. The Southern District of Iowa dismissed, upholding the ban as reasonable and declining to follow the Fifth Circuit's material-and-substantial-interference standard in \textit{Burnside v.\ Byars}. The en banc Eighth Circuit divided equally and affirmed without opinion. The Supreme Court reversed on February 24, 1969.

\section{Judicial Votes}

Reversed and remanded, \textbf{7--2}.
\begin{itemize}
\item \textbf{Majority:} Justice Fortas, joined by Chief Justice Warren and Justices Douglas, Brennan, Stewart, White, and Marshall.
\item \textbf{Concurring:} Justice Stewart; Justice White.
\item \textbf{Dissenting:} Justice Black; Justice Harlan.
\end{itemize}

\section{Holding}

Public schools may not censor passive, non-disruptive student political expression unless officials show facts that might reasonably have led them to forecast ``material and substantial interference'' with schoolwork or discipline, or collision with the rights of other students. Undifferentiated fear or apprehension of disturbance is not enough, and officials may not single out one controversial viewpoint for suppression.

\section{Analysis of Opinions}

\textbf{Majority (Fortas).} Two principles meet: students and teachers keep First Amendment rights in school, as \textit{Barnette} and \textit{Meyer} recognized, yet officials have comprehensive authority to control conduct consistent with constitutional safeguards. The armbands were silent, passive political expression divorced from disruptive conduct, unlike rules on skirt length or hairstyles, and so were ``closely akin to `pure speech.' '' The district court erred in deferring to unsubstantiated fear; any departure from majority opinion may inspire fear, but the Constitution requires that risk, as \textit{Terminiello} held. Adopting \textit{Burnside}, Fortas required more than a desire to avoid the discomfort of an unpopular viewpoint. The record showed no actual or forecast disruption, and banning only antiwar armbands while allowing campaign buttons and Iron Crosses singled out one viewpoint. Schools are not enclaves of totalitarianism, and student intercommunication in hallways and cafeterias is part of education for citizenship.

\textbf{Stewart (concurring).} Stewart rejected any assumption that children's First Amendment rights are coextensive with adults', citing his \textit{Ginsberg v.\ New York} concurrence.

\textbf{White (concurring).} White stressed the distinction between words and expressive conduct that impinges on valid state interests and disclaimed some of \textit{Burnside}'s dicta.

\textbf{Black (dissenting).} Black accused the majority of reviving \textit{Lochner}-style reasonableness review and transferring control of the nation's 23,390 school systems to judges and students. No one may speak wherever and whenever he pleases; the armbands diverted students from their work, including a disrupted mathematics lesson; and the ruling invited a ``new revolutionary era of permissiveness.''

\textbf{Harlan (dissenting).} Officials deserve the widest authority to keep order, and challengers should have to prove a regulation was motivated by viewpoint suppression rather than good-faith educational concerns.

\textbf{Later developments.} \textit{Fraser} (lewd speech), \textit{Hazelwood} (school-sponsored speech), and \textit{Morse} (drug promotion) carved out categories, and \textit{Mahanoy} (2021) sharply limited school authority over off-campus speech.

\section{Examples: Future Applications}

\begin{enumerate}
\item \textbf{Same-Side (Student Wins): A Political T-Shirt.} A student wears a shirt criticizing a school-board climate vote, while ``Save Our Planet'' shirts are allowed. A brief lunchtime argument and the principal's worry about protests are undifferentiated fear, and the selective ban mirrors Des Moines's viewpoint discrimination.
\item \textbf{Same-Side (Student Wins): Lunchtime Leaflets.} A student is suspended for leaflets urging attendance at a city council meeting on police policy. No class was interrupted, and avoiding discomfort is no justification.
\item \textbf{Opposite-Side (School Wins): Documented Disruption.} After fights and walkouts over climate apparel, a school bans all apparel on the dispute, on both sides. Concrete evidence and viewpoint neutrality satisfy \textit{Tinker}.
\item \textbf{Opposite-Side (School Wins): The Class Newspaper.} A principal pulls an article from a school-funded journalism-class paper. \textit{Hazelwood}, not \textit{Tinker}, governs school-sponsored speech.
\item \textbf{Fence-Sitter: A Weekend Post.} A student's off-campus post mocking a teacher circulates in class. \textit{Mahanoy} protects off-campus speech but allows schools to address serious bullying, and \textit{Tinker}'s rights-of-others prong is undefined.
\end{enumerate}

\section{Critique}

\textbf{The heckler's veto.} Because \textit{Tinker} permits censorship when officials reasonably forecast substantial disruption, hostile classmates can arguably silence unpopular speakers by threatening disorder, a tension with \textit{Terminiello} and ordinary hostile-audience doctrine.

\textbf{The rights of others.} Fortas paired material and substantial interference with collision with the rights of other students to be secure and to be let alone, but never defined those rights. Modern courts split sharply over whether psychological discomfort or ideological offense counts as an invasion of peers' rights.

\section{Key Quotations}

\begin{quote}
``It can hardly be argued that either students or teachers shed their constitutional rights to freedom of speech or expression at the schoolhouse gate.''
\end{quote}
\noindent --- Justice Fortas, Opinion of the Court, 393 U.S. at 506.

\begin{quote}
``But, in our system, undifferentiated fear or apprehension of disturbance is not enough to overcome the right to freedom of expression.''
\end{quote}
\noindent --- Justice Fortas, Opinion of the Court, 393 U.S. at 508.

\begin{quote}
``In our system, state-operated schools may not be enclaves of totalitarianism. School officials do not possess absolute authority over their students. \dots\ In our system, students may not be regarded as closed-circuit recipients of only that which the State chooses to communicate.''
\end{quote}
\noindent --- Justice Fortas, Opinion of the Court, 393 U.S. at 511.

\end{document}
